problem:  
Let \(S\) be a standard Einstein solvmanifold in the sense of Heber, with Einstein derivation \(D_E\) and curvature operator \(\mathcal R_E\) under the normalization \(\operatorname{Scal}(g_E) = -1\).  
Denote by \(\Sigma_E = \operatorname{Spec}(\mathcal R_E)\) the (scale‑normalized) spectrum of the curvature operator acting on \(\Lambda^2 T S\), and by \(\sigma_E = \operatorname{Spec}(D_E)\) the spectrum of the Einstein derivation on the Lie algebra level.  

Define the **Einstein–solvable isospectral locus** in the space of \(n\)-dimensional, negatively curved Einstein metrics (modulo diffeomorphism and scaling) by  
\[
\mathcal I_{\Sigma_E,\sigma_E}(n)
  = \left\{ (M,g)\ \middle|\ 
     (M,g)\ \text{Einstein},\ \dim M = n, \operatorname{Scal}(g)=-1,\ 
     \exists\ S' \text{ standard Einstein solvmanifold such that }
     \operatorname{Spec}(\mathcal R_g)=\operatorname{Spec}(\mathcal R_{S'})=\Sigma_E,\ 
     \operatorname{Spec}(D_{S'})=\sigma_E
  \right\}.
\]
This set gathers all Einstein metrics whose **curvature operator spectrum** coincides with that of some standard Einstein solvmanifold having derivation eigenvalues \(\sigma_E\).  

> **Problem (Spectral Rigidity of Einstein Solvmanifolds).**  
> For a fixed dimension \(n\) and a given standard Einstein solvmanifold \(S\), is \(\mathcal I_{\Sigma_E,\sigma_E}(n)\) discrete in the moduli space of Einstein metrics (modulo diffeomorphism and scaling)?  
>  
> Equivalently: does the joint data of curvature‑operator spectrum and derivation spectrum isolate the solvmanifold within the moduli space of all Einstein metrics of the same dimension and scalar normalization, or can there exist continuous families of non‑isometric, non‑homogeneous Einstein metrics sharing exactly the same spectra \((\Sigma_E,\sigma_E)\)?

Why is it a "good" problem:  
This problem asks whether the spectral signatures uniquely identify Einstein solvmanifolds—an analytic–algebraic rigidity question bridging **Riemannian spectral theory**, **Lie algebraic geometry**, and **Einstein moduli theory**.  

1. **Conceptual clarity and depth:** the problem focuses on the discreteness (or lack thereof) of the isospectral locus, a precise and natural notion in the study of moduli spaces; it avoids ill‑defined algebraic constructs beyond standard invariants.  

2. **Bridge between fields:** it links spectral geometry (curvature operator spectra) with Lie algebraic invariants (derivation spectra) and with analytic deformation theory of Einstein metrics. Success would reveal how analytic curvature data reflect or conceal underlying algebraic symmetry—a key theme that further connects to rigidity conjectures of Alekseevskii type.  

3. **Profound simplicity:** the statement is easy to grasp—“can two different Einstein metrics look spectrally identical to a given solvmanifold?”—yet potentially vast in consequence. Understanding it would push forward rigidity frameworks, possibly identifying new invariants separating homogeneous Einstein geometries from nearby non‑homogeneous ones.  

Thus, the problem is mathematically well‑posed, genuinely open, and conceptually rich, satisfying the hallmarks of an excellent research‑level problem.